\documentclass[11pt]{amsart}

\usepackage{amsmath, amssymb, amsthm}
\usepackage{mathtools}
\usepackage{thmtools}
\usepackage{enumitem}
\usepackage{tikz-cd}
\usepackage{hyperref}
\usepackage[nameinlink]{cleveref}

% % Theorem environments
\declaretheorem[numberwithin=section]{theorem}
\declaretheorem[style=plain, sibling=theorem]{lemma, proposition, corollary}
\declaretheorem[style=definition]{definition, example, condition}
\declaretheorem[style=remark]{remark, note, observation, todo}


\title{April 12, 2026}


\begin{document}
\maketitle

Today, I’m unsure how to spend my time. I can either review the paper on intersection theory or work on building a daily system for improving my mathematical writing.
% Today I am struggling for choosing what to do today, to review the paper of intersection theory or to build daily system for self-development for mathematical writing.
\par\vspace{1em}
At first, I talked with ChatGPT about how to configure daily routine for improving mathematical writing skills, and recording research records.
Currently, the routine is temporarily fixed as follows: I write daily research log in a .tex file.
Based on the log = raw file, we perform refinement of the log, and then we also perfom reconstruction of the research result.
These procedures are AI-involved processes.

The principle we concluded is that AI should involve not everytime, it should involve after self-reflection.
In other words, I should first write the raw file by myself, and then I check what I wrote, and then ask AI to evaluate the writing. 

Maybe I need to separate folders for raw files, refined files, and reconstructed files.

\par\vspace{1em}
\begin{todo} To build daily-writing routine,
    \begin{itemize}
        \item Configure Codex skills for evaluating everyday writing.
        It should provide feedback on the writing in the grammatical view point, and the structurual view point.
        \item It should generate a file reviewing for each daily log.
        The resulting file could contain grammatical or English writing feedbacks.
    \end{itemize}    
\end{todo}

The remaining test might be to add another routine for reading papers and studying others research.
This step might include practicing for studying patterns of mathematical writing.
\par\vspace{1em}
We configured the structure of \texttt{raws} folder.
In particular, we added folders for downloaded papers and review files.
Additionally, we added a blueprint file for collecting patterns in mathematical writing.
Currently, written examples are coming from the discussion with ChatGPT.

\begin{todo} There is a todo-list for managing papers.
    \begin{itemize}
        \item Generate a template .tex file for reviewing papers,
        We may just use the template for papers since the review can be regarded as a sort of a math paper.
        \item Configure parsing arxiv bot.
        This bot should download updated papers, which are related to my research topics, in arxiv.
    \end{itemize}
\end{todo}

We write first example paper review file.
The example is based on \cite{Goresky_MacPherson-1980-intersection}.

\bibliographystyle{amsalpha}
% \nocite{*}
\bibliography{bibliography}
\end{document}
